\section{Avaliação e resultados parciais}
\label{sec:resultados}

\subsection{Resultados parciais já obtidos}

A análise documental preliminar produziu quatro resultados verificáveis:

\begin{sloppypar}
\begin{enumerate}
  \item o inventário público do Serpro classifica código e descrição do
  combustível como dados da consulta detalhada de veículo
  \cite{serproDadosVeiculo2026};
  \item a Portaria n\textsuperscript{o} 139/2025 disciplina o acesso aos sistemas e subsistemas
  \cite{senatran1392025}. O fluxo público do Serpro prevê solicitação no
  Credencia, análise da SENATRAN e, após deferimento, contratos com Serpro e GCC
  \cite{serproComoContratar2026};
  \item na Portaria SENATRAN n\textsuperscript{o} 417/2026, a faixa de até 50.000 consultas tem
  valor de R\$~1{,}62 para a modalidade detalhada; R\$~0{,}80 aparece nessa
  modalidade somente na faixa de 2.000.001 a 2.500.000 consultas
  \cite{senatran4172026}; e
  \item os requisitos de segurança para atmosferas explosivas impedem tratar
  componentes eletrônicos comuns da bancada como solução instalável na pista
  \cite{inmetro1152022,nr20}.
\end{enumerate}
\end{sloppypar}

O valor de R\$~1{,}62 é referência pública na data-base de 29 de julho de 2026,
não uma cotação contratual completa. Remuneração da GCC, tributos, infraestrutura,
integração e suporte serão separados na análise econômica. Todos os valores e
normas serão revalidados na data da entrega.

A partir dessas evidências foram elaborados os requisitos das
Tabelas~\ref{tab:rf} e~\ref{tab:rnf}, a arquitetura da
Figura~\ref{fig:arquitetura}, a máquina de estados da Figura~\ref{fig:estados} e
a política da Tabela~\ref{tab:decisao}. Esses elementos são \textbf{artefatos de
projeto preliminares}; ainda não constituem resultado experimental.

Também foi produzido um primeiro incremento do artefato. O backend simulado
implementa consulta de casos sintéticos, mapeamento de combustíveis, decisão,
TTL, intenção persistida antes do comando, ACK síncrono do dispensador em
memória, SQLite minimizado e override restrito a incertezas de
\path{REVIEW_REQUIRED}. A idempotência por \path{client_request_id}, o
fingerprint HMAC persistido e o cache HMAC temporário evitam, nos casos cobertos,
novo comando ou segunda consulta durante a reconciliação. A PWA implementa o
fluxo operacional em modo demonstrativo ou conectado à API e trata resposta
perdida como estado desconhecido até reconciliação manual. O firmware
experimental implementa a lógica local de TTL, sequência e idempotência.

Em execução local de 29 de julho de 2026, os 82 testes automatizados do backend
foram aprovados e a ferramenta de cobertura reportou 93\% das instruções. Os
12 testes da PWA também foram aprovados, sua compilação de produção foi
concluída, o executável de teste da lógica C++ do firmware terminou sem falha e
o alvo \path{esp32dev} foi compilado com PlatformIO.
Esse resultado é reprodutível pelos comandos documentados no repositório, mas
se limita a componentes de software. Não houve teste ponta a ponta entre PWA,
backend, broker e ESP32, gravação/execução em placa nem ensaio com bancada física.

\subsection{Resultados ainda não disponíveis}

\begin{table}[H]
  \centering
  \caption{Evidências necessárias para testar as hipóteses}
  \label{tab:resultados-pendentes}
  \small
  \begin{tabularx}{\textwidth}{@{}L{3.2cm}YL{2.5cm}@{}}
    \toprule
    \textbf{Evidência} & \textbf{Uso} & \textbf{Estado} \\
    \midrule
    Ocorrências e custos anonimizados & Caracterizar o caso e alimentar $L$ & Não coletada \\
    Total de abastecimentos compatível & Calcular taxa, se houver cobertura equivalente & Não coletado \\
    Entrevistas autorizadas & Refinar processo, requisitos e modos de falha & Não realizadas \\
    Matriz integral de H1 & Testar decisões, falhas e bloqueio em todos os casos predefinidos & Parcial: testes de componentes \\
    Integração MQTT/ESP32 & Confirmar comando, sessão, TTL e retorno ao bloqueio & Não executada \\
    Ciclos da bancada & Testar repetibilidade e intertravamento físico & Não executados \\
    Medidas de latência/usabilidade & Avaliar desempenho e interação & Não coletadas \\
    Custos da solução & Calcular $C_{\mathrm{sol}}$ & Incompletos \\
    \bottomrule
  \end{tabularx}
  \\[2pt]\fontepropria
\end{table}

Assim, os testes existentes são evidência preliminar de construção, mas H1 e H2
permanecem sem teste conclusivo. Esta versão não afirma redução de incidentes,
retorno financeiro ou adequação para operação real.

\subsection{Critérios de aceitação}

H1 exigirá zero liberação incompatível em todos os casos predefinidos, zero
liberação automática diante de resposta inválida, ambígua ou indisponível e
retorno ao bloqueio após expiração, término ou falha. Serão incluídos gasolina,
etanol, flex, diesel S10, diesel S500, diesel sem especificação inequívoca,
híbrido, elétrico, conversão não comprovada, veículo não encontrado e provedor
indisponível. A quantidade de repetições será definida com
o orientador antes dos ensaios e não será escolhida depois de observar os
resultados.

H2 será avaliada nos cenários de 25\%, 50\%, 75\% e 100\% de prevenção. A
hipótese será suportada se $B(p)>0$ em pelo menos um cenário, desde que $L$ e
$C_{\mathrm{sol}}$ tenham dados adequados. Caso os registros sejam insuficientes, a
análise será apresentada como exercício de sensibilidade e H2 ficará
inconclusiva, não rejeitada por conveniência.
